\documentclass{article}
\usepackage{amsmath,amssymb,amsthm}
\usepackage{enumitem}
\newtheorem{theorem}{Theorem}
\newtheorem{lemma}{Lemma}
\newtheorem{assumption}{Assumption}

\begin{document}

\section*{Algorithm Name}
\textbf{Stochastic Gradient Descent with Decreasing Noise (SGD‑DN)}  

The method performs a standard gradient step on a convex smooth objective and adds a zero‑mean random perturbation whose variance decays with the iteration index.

\section*{Mathematical Formulation}
Let $f:\mathbb{R}^{d}\rightarrow\mathbb{R}$ be the objective function and denote by 
\[
x_{k}\in\mathbb{R}^{d},\qquad k=0,1,2,\dots
\]
the iterate at iteration $k$.  The update rule is
\begin{equation}
\boxed{ \; x_{k+1}=x_{k}-\alpha\,\nabla f(x_{k})+\xi_{k}\; } \tag{1}
\end{equation}
where
\begin{itemize}
    \item $\alpha>0$ is a fixed step size (learning rate);
    \item $\xi_{k}\in\mathbb{R}^{d}$ is a random noise vector satisfying the conditions listed in the \textbf{Assumptions} section below.
\end{itemize}
The noise variance is allowed to shrink, e.g.
\begin{equation}
\mathbb{E}\!\left[\|\xi_{k}\|^{2}\mid x_{k}\right]\le \frac{\sigma^{2}}{k+1},\qquad \sigma>0. \tag{2}
\end{equation}

\section*{Pseudocode}
\begin{enumerate}[label=\textbf{Step \arabic*:},leftmargin=*]
\item \textbf{Input:} step size $\alpha$, initial point $x_{0}$, total iterations $K$, noise schedule $\{\sigma_{k}^{2}\}_{k\ge0}$.
\item \textbf{For} $k=0,1,\dots ,K-1$ \textbf{do}
    \begin{enumerate}[label=\arabic*.]
        \item Compute the (exact) gradient $g_{k}\gets\nabla f(x_{k})$.
        \item Draw a random vector $\xi_{k}$ with $\mathbb{E}[\xi_{k}\mid x_{k}]=0$ and 
        $\mathbb{E}[\|\xi_{k}\|^{2}\mid x_{k}]\le\sigma_{k}^{2}$.
        \item Update $x_{k+1}\gets x_{k}-\alpha g_{k}+\xi_{k}$.
    \end{enumerate}
\item \textbf{Output:} $x_{K}$ (or the average $\bar x_{K}=K^{-1}\sum_{k=0}^{K-1}x_{k}$).
\end{enumerate}

\section*{Dry Run Example}
Consider a one–dimensional quadratic
\[
f(w)=(w-3)^{2},
\]
so that $\nabla f(w)=2(w-3)$.  Choose $\alpha=0.1$, $w_{0}=0$, and a deterministic “noise’’ sequence $\xi_{k}=0$ (to illustrate the pure deterministic part).  

\begin{center}
\begin{tabular}{c|c|c|c}
$k$ & $w_{k}$ & $\nabla f(w_{k})$ & $w_{k+1}=w_{k}-\alpha\nabla f(w_{k})$ \\ \hline
0 & $0$ & $-6$ & $0-0.1(-6)=0.6$ \\
1 & $0.6$ & $-4.8$ & $0.6-0.1(-4.8)=1.08$ \\
2 & $1.08$ & $-3.84$ & $1.08-0.1(-3.84)=1.464$ \\
\end{tabular}
\end{center}
The corresponding objective values are $f(w_{0})=9$, $f(w_{1})=5.76$, $f(w_{2})=3.6864$, showing a steady decrease.

\newpage

\section*{Assumptions}
\begin{assumption}[Convexity]\label{ass:convex}
$f$ is convex, i.e. for all $x,y\in\mathbb{R}^{d}$
\[
f(y)\ge f(x)+\langle\nabla f(x),y-x\rangle .
\]
\end{assumption}
\begin{assumption}[L‑smoothness]\label{ass:smooth}
$f$ has $L$‑Lipschitz continuous gradient:
\[
\|\nabla f(x)-\nabla f(y)\|\le L\|x-y\|,\qquad\forall x,y\in\mathbb{R}^{d}.
\]
Equivalently,
\[
f(y)\le f(x)+\langle\nabla f(x),y-x\rangle+\frac{L}{2}\|y-x\|^{2}.
\]
\end{assumption}
\begin{assumption}[Noise]\label{ass:noise}
For each iteration $k$ the random vector $\xi_{k}$ satisfies
\[
\mathbb{E}[\xi_{k}\mid x_{k}]=0,\qquad 
\mathbb{E}\!\left[\|\xi_{k}\|^{2}\mid x_{k}\right]\le \frac{\sigma^{2}}{k+1},
\]
with a constant $\sigma>0$.
\end{assumption}
\begin{assumption}[Step size]\label{ass:stepsize}
The step size $\alpha$ is constant and chosen such that
\[
0<\alpha\le\frac{1}{L}.
\]
\end{assumption}

\section*{Main Convergence Theorem}
\begin{theorem}[Expected suboptimality of SGD‑DN]\label{thm:main}
Let $x^{\star}\in\arg\min_{x}f(x)$ and assume \ref{ass:convex}–\ref{ass:stepsize}.  Define the averaged iterate
\[
\bar x_{K}\;:=\;\frac{1}{K}\sum_{k=0}^{K-1}x_{k}.
\]
Then for any $K\ge1$,
\[
\frac{1}{K}\sum_{k=0}^{K-1}\mathbb{E}\!\bigl[f(x_{k})-f(x^{\star})\bigr]
\;\le\;
\frac{\|x_{0}-x^{\star}\|^{2}}{2\alpha K}
\;+\;
\frac{\alpha\sigma^{2}}{K}\sum_{k=0}^{K-1}\frac{1}{k+1}
\;\le\;
\frac{\|x_{0}-x^{\star}\|^{2}}{2\alpha K}
\;+\;
\frac{\alpha\sigma^{2}\bigl(\log K+1\bigr)}{K}.
\tag{3}
\]
Consequently,
\[
\mathbb{E}\!\bigl[f(\bar x_{K})-f(x^{\star})\bigr]=O\!\left(\frac{1}{K}\right)
\quad\text{and}\quad
\mathbb{E}\!\bigl[f(x_{K})-f(x^{\star})\bigr]=O\!\left(\frac{\log K}{K}\right).
\]
\end{theorem}

\section*{Theoretical Properties}
\begin{enumerate}[label=\arabic*.]
\item \textbf{Stability.} The requirement $\alpha\le1/L$ guarantees that the deterministic part of the recursion is non‑expansive (see Lemma~\ref{lem:one_step}).
\item \textbf{Hyper‑parameter sensitivity.} The bound (3) shows a linear dependence on $\alpha$ in the variance term; a smaller $\alpha$ reduces the steady‑state error caused by noise but slows down the $1/(\alpha K)$ term.
\item \textbf{Comparison to baselines.}
    \begin{itemize}
        \item \textit{Gradient descent (GD)} corresponds to $\xi_{k}=0$ and recovers the classic $O(1/K)$ rate without the extra variance term.
        \item \textit{Standard SGD} with i.i.d.\ bounded variance $\mathbb{E}\!\left[\|\xi_{k}\|^{2}\right]\le\sigma^{2}$ yields a bound of order $O(1/\sqrt{K})$; the decreasing variance schedule improves the rate to $O(1/K)$ (up to a logarithmic factor).
    \end{itemize}
\item \textbf{Asymptotic behavior.} Since $\sum_{k=0}^{\infty}\sigma^{2}/(k+1)=\infty$, the variance term cannot be eliminated completely, but its contribution vanishes as $K\to\infty$ because of the $1/K$ prefactor.
\end{enumerate}

\section*{Discussion}
The theorem demonstrates that, under convexity and smoothness, adding a *vanishing* stochastic perturbation does not destroy the $O(1/K)$ convergence of deterministic gradient descent.  The proof mirrors the classic SGD analysis, with the only modification being the explicit decay $\sigma^{2}/(k+1)$ in the noise variance, which yields a harmonic series that contributes only a $\log K$ factor.

\newpage

\section*{Preliminaries (Definitions and Lemmas)}
\begin{lemma}[One‑step inequality]\label{lem:one_step}
Under Assumptions \ref{ass:smooth}–\ref{ass:stepsize} and the update (1),
\[
\|x_{k+1}-x^{\star}\|^{2}
\le
\|x_{k}-x^{\star}\|^{2}
-2\alpha\bigl(f(x_{k})-f(x^{\star})\bigr)
+\alpha^{2}\|\xi_{k}\|^{2}.
\tag{4}
\]
\end{lemma}
\begin{proof}
\begin{enumerate}[label=[Step \arabic*],leftmargin=*]
\item Expand the squared norm using the definition of $x_{k+1}$:
\[
\|x_{k+1}-x^{\star}\|^{2}
= \|x_{k}-\alpha\nabla f(x_{k})+\xi_{k}-x^{\star}\|^{2}.
\tag{5}
\]
\item Apply $\|a+b\|^{2}= \|a\|^{2}+2\langle a,b\rangle+\|b\|^{2}$ with 
$a:=x_{k}-\alpha\nabla f(x_{k})-x^{\star}$ and $b:=\xi_{k}$:
\[
\|x_{k+1}-x^{\star}\|^{2}
= \|x_{k}-\alpha\nabla f(x_{k})-x^{\star}\|^{2}
+2\langle x_{k}-\alpha\nabla f(x_{k})-x^{\star},\xi_{k}\rangle
+\|\xi_{k}\|^{2}.
\tag{6}
\]
\item Take conditional expectation w.r.t.\ $\mathcal{F}_{k}:=\sigma(x_{0},\dots ,x_{k})$.
Because $\mathbb{E}[\xi_{k}\mid\mathcal{F}_{k}]=0$ (Assumption~\ref{ass:noise}),
the cross term vanishes:
\[
\mathbb{E}\!\bigl[\langle x_{k}-\alpha\nabla f(x_{k})-x^{\star},\xi_{k}\rangle\mid\mathcal{F}_{k}\bigr]=0.
\tag{7}
\]
\item Hence
\[
\mathbb{E}\!\bigl[\|x_{k+1}-x^{\star}\|^{2}\mid\mathcal{F}_{k}\bigr]
=
\|x_{k}-\alpha\nabla f(x_{k})-x^{\star}\|^{2}
+\mathbb{E}\!\bigl[\|\xi_{k}\|^{2}\mid\mathcal{F}_{k}\bigr].
\tag{8}
\]
\item Expand the deterministic part:
\[
\|x_{k}-\alpha\nabla f(x_{k})-x^{\star}\|^{2}
=
\|x_{k}-x^{\star}\|^{2}
-2\alpha\langle\nabla f(x_{k}),x_{k}-x^{\star}\rangle
+\alpha^{2}\|\nabla f(x_{k})\|^{2}.
\tag{9}
\]
\item By convexity (Assumption~\ref{ass:convex}),
\[
f(x_{k})-f(x^{\star})\le\langle\nabla f(x_{k}),x_{k}-x^{\star}\rangle .
\tag{10}
\]
\item By $L$‑smoothness (Assumption~\ref{ass:smooth}) and the step‑size bound $\alpha\le1/L$,
\[
\alpha^{2}\|\nabla f(x_{k})\|^{2}
\le 2\alpha\bigl(f(x_{k})-f(x^{\star})\bigr).
\tag{11}
\]
(Indeed, $ \|\nabla f(x_{k})\|^{2}\le 2L\bigl(f(x_{k})-f(x^{\star})\bigr)$, multiply by $\alpha^{2}\le\alpha/L$.)

\item Substitute (10) and (11) into (9):
\[
\|x_{k}-\alpha\nabla f(x_{k})-x^{\star}\|^{2}
\le
\|x_{k}-x^{\star}\|^{2}
-2\alpha\bigl(f(x_{k})-f(x^{\star})\bigr)
+2\alpha\bigl(f(x_{k})-f(x^{\star})\bigr)
=
\|x_{k}-x^{\star}\|^{2}
-2\alpha\bigl(f(x_{k})-f(x^{\star})\bigr).
\tag{12}
\]
\item Combine (8) and (12) and drop the conditional expectation notation for brevity,
yielding exactly (4).  ∎
\end{enumerate}
\end{proof}

\section*{Main Proof (Step‑by‑Step Derivation)}
We now prove Theorem~\ref{thm:main}.

\begin{enumerate}[label=[Step \arabic*],leftmargin=*]
\item Take total expectation of Lemma~\ref{lem:one_step} (4):
\[
\mathbb{E}\!\bigl[\|x_{k+1}-x^{\star}\|^{2}\bigr]
\le
\mathbb{E}\!\bigl[\|x_{k}-x^{\star}\|^{2}\bigr]
-2\alpha\,\mathbb{E}\!\bigl[f(x_{k})-f(x^{\star})\bigr]
+\alpha^{2}\mathbb{E}\!\bigl[\|\xi_{k}\|^{2}\bigr].
\tag{13}
\]
\item Apply the variance bound (2):
\[
\mathbb{E}\!\bigl[\|\xi_{k}\|^{2}\bigr]\le\frac{\sigma^{2}}{k+1}.
\tag{14}
\]
\item Substitute (14) into (13):
\[
\mathbb{E}\!\bigl[\|x_{k+1}-x^{\star}\|^{2}\bigr]
\le
\mathbb{E}\!\bigl[\|x_{k}-x^{\star}\|^{2}\bigr]
-2\alpha\,\mathbb{E}\!\bigl[f(x_{k})-f(x^{\star})\bigr]
+\frac{\alpha^{2}\sigma^{2}}{k+1}.
\tag{15}
\]
\item Rearrange (15) to isolate the suboptimality term:
\[
2\alpha\,\mathbb{E}\!\bigl[f(x_{k})-f(x^{\star})\bigr]
\le
\mathbb{E}\!\bigl[\|x_{k}-x^{\star}\|^{2}\bigr]
-
\mathbb{E}\!\bigl[\|x_{k+1}-x^{\star}\|^{2}\bigr]
+\frac{\alpha^{2}\sigma^{2}}{k+1}.
\tag{16}
\]
\item Sum (16) over $k=0,\dots ,K-1$:
\[
2\alpha\sum_{k=0}^{K-1}\mathbb{E}\!\bigl[f(x_{k})-f(x^{\star})\bigr]
\le
\|x_{0}-x^{\star}\|^{2}
-
\mathbb{E}\!\bigl[\|x_{K}-x^{\star}\|^{2}\bigr]
+\alpha^{2}\sigma^{2}\sum_{k=0}^{K-1}\frac{1}{k+1}.
\tag{17}
\]
(The telescoping sum collapses the first two terms.)
\item Drop the non‑negative term $-\mathbb{E}[\|x_{K}-x^{\star}\|^{2}]$ to obtain an upper bound:
\[
2\alpha\sum_{k=0}^{K-1}\mathbb{E}\!\bigl[f(x_{k})-f(x^{\star})\bigr]
\le
\|x_{0}-x^{\star}\|^{2}
+\alpha^{2}\sigma^{2}\sum_{k=0}^{K-1}\frac{1}{k+1}.
\tag{18}
\]
\item Divide both sides by $2\alpha K$:
\[
\frac{1}{K}\sum_{k=0}^{K-1}\mathbb{E}\!\bigl[f(x_{k})-f(x^{\star})\bigr]
\le
\frac{\|x_{0}-x^{\star}\|^{2}}{2\alpha K}
+\frac{\alpha\sigma^{2}}{K}\sum_{k=0}^{K-1}\frac{1}{k+1}.
\tag{19}
\]
\item The harmonic sum satisfies $\sum_{k=0}^{K-1}\frac{1}{k+1}\le\log K+1$.  Plugging this bound yields (3).
\end{enumerate}
Thus the claimed rate follows.  Moreover, by Jensen’s inequality,
\[
\mathbb{E}\!\bigl[f(\bar x_{K})-f(x^{\star})\bigr]
\le
\frac{1}{K}\sum_{k=0}^{K-1}\mathbb{E}\!\bigl[f(x_{k})-f(x^{\star})\bigr],
\]
so the same $O(1/K)$ bound holds for the averaged iterate as well.

\section*{Rate Derivation (With Constants)}
From (19) we can read off the two contributions:
\[
\underbrace{\frac{\|x_{0}-x^{\star}\|^{2}}{2\alpha K}}_{\text{deterministic GD term}}
\qquad\text{and}\qquad
\underbrace{\frac{\alpha\sigma^{2}(\log K+1)}{K}}_{\text{variance term}}.
\]
Choosing $\alpha=1/L$ (the largest admissible step) balances the two terms most favorably:
\[
\mathbb{E}\!\bigl[f(\bar x_{K})-f(x^{\star})\bigr]
\le
\frac{L\|x_{0}-x^{\star}\|^{2}}{2K}
+\frac{\sigma^{2}(\log K+1)}{L K}.
\]
Hence the dominant asymptotic order is $O(1/K)$.

\section*{Special Cases}
\begin{enumerate}[label=\alph*.]
\item \textbf{Zero noise ($\sigma=0$).}  The bound reduces to the classic deterministic gradient descent rate $\|x_{0}-x^{\star}\|^{2}/(2\alpha K)$.
\item \textbf{Constant variance ($\sigma_{k}^{2}\equiv\sigma^{2}$).}  The same derivation would give a term $\alpha\sigma^{2}$ (no $1/k$ factor), leading to an $O(1/\sqrt{K})$ rate after optimizing $\alpha$, matching standard SGD results.
\item \textbf{Strongly convex $f$ (parameter $\mu>0$).}  Adding the inequality 
$\langle\nabla f(x_{k}),x_{k}-x^{\star}\rangle\ge \mu\|x_{k}-x^{\star}\|^{2}$
to Lemma~\ref{lem:one_step} yields a linear convergence factor $(1-2\alpha\mu)$ plus a vanishing variance floor, i.e.
\[
\mathbb{E}\!\bigl[f(x_{k})-f(x^{\star})\bigr]
\le
(1-2\alpha\mu)^{k}\bigl(f(x_{0})-f(x^{\star})\bigr)
+\frac{\alpha\sigma^{2}}{2\mu(k+1)}.
\]
\end{enumerate}

\end{document}